\chapter{Appendices}
\label{chap:annexes}

\section{Keyboard shortcuts}

\begin{tabularx}{\linewidth}{l X}
    \toprule
    \textbf{Shortcut} & \textbf{Action} \\
    \midrule
    \touche{Ctrl}+\touche{O}            & Open current airfoil \\
    \touche{Ctrl}+\touche{Shift}+\touche{O} & Open reference airfoil \\
    \touche{Ctrl}+\touche{S}            & Save current airfoil \\
    \touche{Ctrl}+\touche{Q}            & Quit \\
    \touche{Ctrl}+\touche{Z}            & Undo (reserved, not active) \\
    \touche{Ctrl}+\touche{Y}            & Redo (reserved, not active) \\
    \touche{Ctrl}+\touche{B}            & Convert to spline \\
    \touche{Ctrl}+\touche{0}            & Fit zoom (airfoils canvas) \\
    \bottomrule
\end{tabularx}

\section{Glossary}

\begin{description}
    \item[LE / TE] Leading edge ($x=0$) / Trailing edge ($x=c$).
    \item[Camber] Maximum height of the airfoil mean line,
        expressed as a percentage of the chord.
    \item[$C_D$] Drag coefficient. Drag force
        normalized by the dynamic pressure and the chord.
    \item[$C_f$] Local wall skin-friction coefficient.
    \item[$C_L$] Lift coefficient.
    \item[$C_M$] Pitching moment coefficient, generally
        referenced to the quarter chord.
    \item[$C_p$] Pressure coefficient: $(p - p_\infty) /
        (\frac{1}{2}\rho V_\infty^2)$.
    \item[$C^0$, $C^1$, $C^2$ continuity] At the junction point of two curves:
        identity of position ($C^0$), of tangent
        ($C^1$), of curvature ($C^2$).
    \item[Chord] Distance between the LE and the TE, used as the
        reference length.
    \item[Bézier curve] Parametric curve defined by a
        control polygon. The degree equals the number of
        control points minus one.
    \item[$\delta^*$] Boundary-layer displacement thickness.
    \item[De Casteljau] Recursive algorithm for evaluating and
        subdividing a Bézier curve. The subdivision is
        \emph{exact} (no approximation).
    \item[Stall] Abrupt loss of lift beyond a
        critical angle of attack, due to boundary-layer
        separation.
    \item[Upper surface / Lower surface] Upper / lower surface of the
        airfoil.
    \item[L/D] $C_L/C_D$ ratio. The higher the lift-to-drag ratio,
        the better the aerodynamic efficiency.
    \item[$H$] Boundary-layer shape factor,
        $H = \delta^*/\theta$.
    \item[NCRIT] Transition criterion in the $e^N$ method;
        maximum accepted amplitude of disturbances before
        the onset of turbulence.
    \item[Polar] Plot of one aerodynamic quantity as a
        function of another, at a given Reynolds number. Several
        polars form a family indexed by the Reynolds number.
    \item[Reynolds] Dimensionless number
        $\mathrm{Re} = \rho V c / \mu$, characterizing the flow
        regime (laminar, transitional, turbulent).
    \item[Selig] Convention for storing the points of an airfoil:
        TE $\to$ upper surface $\to$ LE $\to$ lower surface $\to$ TE.
    \item[Spline] Piecewise curve composed of several
        Bézier segments joined according to a given level of
        continuity.
    \item[$\theta$] Boundary-layer momentum thickness.
    \item[Trip / Turbulator] Element (roughness, stretched wire) forcing
        the laminar-to-turbulent transition at a given position.
    \item[$U_e$] Velocity at the outer edge of the boundary layer.
    \item[XFoil] 2D aerodynamic solver developed by Mark Drela
        (MIT) coupling a panel method with boundary-layer
        integral equations.
    \item[XTR] Boundary-layer transition position ($x_{tr}/c$).
\end{description}

\section{Troubleshooting (FAQ)}

\subsection*{The GUI does not launch}

Check that:

\begin{enumerate}
    \item the virtual environment is activated (\chemin{env\_py3});
    \item the dependencies are installed (\code{pip install -r
        requirements.txt});
    \item Python is version 3.10 or higher
        (\code{python --version}).
\end{enumerate}

\subsection*{XFoil does not run}

\begin{itemize}
    \item Check for the presence of \chemin{externaltools/xfoil/xfoil.exe}.
    \item On Windows, some antivirus programs may block
        execution. Add an exception for the
        \chemin{externaltools} folder.
    \item The commands sent to XFoil are visible in the
        console (in development mode); review the output
        to diagnose any syntax error.
\end{itemize}

\subsection*{XFoil does not converge for some angles}

\begin{itemize}
    \item Disable \menu{ASEQ} to switch to bidirectional ALFA
        mode.
    \item Reduce the $\alpha$ step (for example from 0.5\textdegree{} to
        0.25\textdegree{}) to ease continuation tracking.
    \item Increase \menu{Timeout} if the computation is long.
    \item Check that the airfoil has no geometric
        singularities (duplicated points, upper/lower surface crossing,
        degenerate LE).
\end{itemize}

\subsection*{The curvature is not displayed}

\begin{itemize}
    \item Convert the airfoil to spline mode
        (\menu{Edit $\rightarrow$ Convert to Spline}).
    \item Increase the porcupine scale (right-click on the
        canvas $\rightarrow$ \menu{Porcupine scale\dots}) if
        they are too short to be visible.
\end{itemize}

\subsection*{I lose my splines when saving}

The \chemin{.dat} format (Selig/Lednicer) does not preserve the
spline representation: only the discrete points are
saved. To keep the spline, use the
\chemin{.bspl} format.

\section{Further reading}

\begin{itemize}
    \item \textbf{Technical documentation}: see the docstrings
        in the source code (Sphinx reST format) and the documentation
        generated in the \chemin{docs/} folder.
    \item \textbf{Unit tests}: see the
        \chemin{tests/} folder, which contains more than 240 tests
        (Bezier, Profil, BezierSpline, ProfilSpline,
        Simulation).
    \item \textbf{Reference article on XFoil}:
        Drela, M. \emph{XFOIL: An Analysis and Design System for
        Low Reynolds Number Airfoils}, in \emph{Low Reynolds
        Number Aerodynamics}, Springer Lecture Notes in Engineering,
        Vol.~54, 1989.
    \item \textbf{Airfoil library}:
        \url{https://m-selig.ae.illinois.edu/ads/coord_database.html}
        (large catalogue in \chemin{.dat} format, accessible
        directly from the application's \menu{File $\rightarrow$ From
        the UIUC database} menu).

\end{itemize}

\bigskip
\hrule
\bigskip

\begin{center}
    \itshape
    End of the user manual.\\
    To report a typo, suggest an improvement or
    ask a question, open an \emph{issue} on
    \chemin{https://github.com/BenoitGFree/AirfoilTools}.
\end{center}
